\section{几何自编码器：关键技术挑战}

\subsection{问题：图像 VAE 不适用于几何数据}

B3D 使用隐扩散模型同时生成 RGB 图像和 XYZ 点图，这要求对两种模态都进行自编码。然而，Henzler 发现\textbf{直接使用预训练的图像 VAE 来编解码点图会产生极差的结果}，且场景越大效果越差。

原因分析：
\begin{itemize}
    \item \textbf{损失函数不匹配}：图像压缩常用的感知损失（perceptual loss）对几何数据毫无意义
    \item \textbf{数据范围不匹配}：图像像素值始终在 $[0, 1]$ 范围内，但几何坐标的取值范围完全不同，且随场景规模变化
\end{itemize}

\begin{warningbox}{图像 VAE 用于几何编码的陷阱}
将为图像设计的 VAE（使用感知损失、GAN 损失等训练）直接应用于几何数据（XYZ 点图）是行不通的。几何数据与图像数据在统计分布、值域范围和语义含义上都有本质区别。必须设计专门的几何自编码器。
\end{warningbox}

\subsection{解码器架构：Transformer vs. 卷积}

另一个有趣的发现是关于解码器架构的选择：

\begin{itemize}
    \item 从定量指标看，卷积解码器和 Transformer 解码器的几何重建质量相当
    \item 但\textbf{定性观察}揭示了关键差异：卷积解码器会产生感知上令人不适的伪影，例如无法保持直线的平直性
    \item 因此 B3D 选择了 \textbf{Transformer 解码器}来处理 XYZ 数据
\end{itemize}

\begin{warningbox}{定量指标 vs. 定性评估的差异}
这是一个典型案例：定量指标（如 MSE）可能无法完全反映实际的视觉质量。两种解码器在数值指标上表现相似，但卷积解码器的输出存在结构性伪影（如弯曲的直线），这在实际应用中是不可接受的。这提醒我们在评估几何重建质量时，不能仅依赖传统指标。
\end{warningbox}

\subsection{本章小结}

几何自编码器是3D隐扩散模型的关键组件，但其设计面临与图像自编码器截然不同的挑战。数据分布差异、损失函数选择和解码器架构都需要针对几何数据特点重新考量。

